\chapter{Minimal Sufficiency}

A sufficient statistc is minial if it is a function of any other sufficient statistics

\subsection{Certification of minimality}

The minimality of a sufficient statistic can be verified based on the following proposition: a sufficient statistic $T$ is minimal if for any $x, y \in \mathcal{X}$,
\[
    \frac{f_\theta(x)}{f_\theta(y)} \text{is independent of } \theta \implies T(x) = T(y)
\]
if the ratio of their likelihoods is independent of $\theta$, then need to check if it implies $T(x) = T(y)$. Only a sufficient condition.

Note that
\[
    \frac{f_\theta(x)}{f_\theta(y)} = \left( \frac{\theta}{1-\theta}\right)^{T(x)-T(y)}
\]
Clearly, the ratio does not depend on $\theta$ iff $T(x)=T(y)$.


\marginpar{\small\textbf{Jan. 20}}

Any function of the observation that preserves all the information is a sufficient statistic.

\subsubsection{Minimal sufficient statistic}

A sufficient statistic is \textit{minimal} if it is a function of any other sufficient statistics.
This is the "optimal" data compression.

A sufficient condition for minimality given sufficiency:
The minimality of a sufficient statistic can be verified based on the following proposition:
A sufficient statistic $T$ is minimal for any $x, y \in \mathcal{X}$,
\[
    \frac{f_\theta(x)}{f_\theta(y)} \text{ is independent of} \theta \implies T(x) = T(y)
\]

(Proof) Let $T$ be a sufficient statistic that staisfies the above property. To show that $T$
is minimal, we must show that for any other statistic $T'$, $T$ is a function of $T'$.

Now suppose that $T'(x) = T'(y)$. We have
\[
    \frac{f_\theta(x)}{f_\theta(y)} = \frac{g_\theta'(T'(x))h'(x)}{g_\theta'(T'(y))h'(y)} = \frac{h'(x)}{h'(y)}
\]
which is independent of $\theta$

Previously, we assume $T$ is a sufficient statistic and then provides a sufficient confition for the minimality.

Let $T=T(X)$ be a stat. It can be shown that $T$ is a sufficient statistic if for any $x, y \in \mathcal{X}$,
\[
    T(x)=T(y) \implies \frac{f_\theta(x)}{f_\theta(y)} \text{ is independent of } \theta.
\]
Proof based on Fisher-Neyman...

Combined with the previous proposition, a statistic $T=T(X)$ is a minimal sufficient statistic
if for any $x, y, \in \mathcal{X}$,
\[
    \frac{f_\theta(x)}{f_\theta(y)} \text{is independent of} \theta \leftrightarrow T(x) = T(y)
\]

\begin{exmp}
    [Bernoulli trials]
    Consider the statistic
    \[
        T = \sum_{i=1}^{n} X_i
    \]
    Note that
    \[
        \frac{f_\theta(x)}{f_\theta(y)} = \left(\frac{\theta}{1-\theta}\right)^{T(x)-T(y)}
    \]
    Clearly, the ratio $\frac{f_\theta(x)}{f_\theta(y)}$ does not dpeend on $\theta$ iff $T(x) = T(y)$.
\end{exmp}

\subsubsection{Complete sufficient statistic}

A statistic $T=T(X)$ is complete if for any real-valued function $\phi$,
\[
    \mathbb{E}_\theta [\phi(T)] = 0, \forall \theta \in \Theta \implies \mathbb{P}_\theta(\phi(T)=0) = 0, \forall 0 \in \Theta
\]


\begin{exmp}
    [Bernoullie trials]
    Recall
    \[
        T = \sum_{i=1}^{n} X_i
    \]
    is a sufficient statistic for $\theta$.

    If for some real-valued function $\phi$ we have
    \[
        \mathbb{E}_\theta[\phi(T)] = \sum_{t=0}^{n} \phi(t) {n \choose t} \theta^t (1-\theta)^{n-t} = 0
    \]
    we must have
    \[
        \sum_{t=0}^{n} \phi(t) {n \choose t} \left( \frac{\theta}{1-\theta} \right)^t = 0
    \]
    for all $\theta \in (0, 1)$.

\end{exmp}

\begin{exmp}
    [Uniform with unknwon support]
    Recall that for uniform random variables over an unknown support $[0, \theta]$.
    \[
        T = \max_{i \in [n]} X_i
    \]
    is a sufficient statistic for $\theta$.
    The pdf of $T$ can be calculated as
    \[
        f_\theta(t) = n t^{n-1} \theta^{-n} 1_{0 < t \leq \theta}
    \]
    If for some real-valued function $\phi$ we have
    \[
        \mathbb{E}_\theta[\phi(T)] = n \theta^{-n} \int_0^\theta \phi(t) t^{n-1} dt = 0
    \]
    for all $\theta > 0$, we must have
    \[
        g(\theta) = \int_0^\theta \phi(t) t^{n-1} dt = 0, \quad \forall \theta > 0
    \]
    Taking the derivative, $g'(\theta)=\phi(\theta)\theta^{n-1}$.

\end{exmp}]

\marginpar{\small\textbf{Jan. 22}}

\section{Complete sufficient statistic}
Under mild conditions, it can be shown that if a sufficient statistic is complete, it is also minimal. \\

Let $T$ be a complete sufficient statistic and $T'$ be another sufficient statistic. We need to show that $T$ is a function of $T'$. \\

Let
\begin{align*}
    \psi(t') & := \mathbb{E}_\theta [T | T' = t']   \\
    \rho(t)  & := \mathbb{E}_\theta[\psi(T') | T=t]
\end{align*}
By law of total expectation,
\begin{align*}
    \mathbb{E}_\theta & = \mathbb{E}_\theta[\mathbb{E}_\theta[T|T']]       \\
                      & = \mathbb{E}_\theta[\psi(T')]                      \\
                      & = \mathbb{E}_\theta[\mathbb{E}_\theta[\psi(T')|T]] \\
                      & = \mathbb{E}_\theta[\rho(T)]
\end{align*}
for any $\theta \in \Theta$. By the completeness of $T$, we have $\mathbb{P}_\theta(T=\rho(T)) = 1$. \\

\noindent Recall the law of total variance: $\Var[Y] = \E[\Var[Y|Z]]+\Var[\E[Y|Z]]$. Now we have
\begin{align*}
    \Var_\theta[\rho_i(T)] & = \E_\theta[\Var_\theta[\rho_i(T)|T']] + \E_\theta[\Var_\theta[\psi_i(T')|T]] + ...
\end{align*}